\documentclass[a4paper]{article}     
\usepackage{amsfonts}
\usepackage{amsmath}
\usepackage{amssymb}
\usepackage{graphicx}

\begin{document}

\noindent \large Auteur: Daan Kallijarvi \\
\noindent \large Studierichting: Informatica\\
\noindent \large Oefeningenreeks 2D nr. 20.10\\

\medskip

\normalsize

Bereken $\displaystyle\lim_{x \to \pm\infty} \dfrac{2 - |x|}{3x + 3|x|}$.\\

\textbf{Geval 1: $x \to -\infty$}\\

Voor $x < 0$ geldt $|x| = -x$, dus de noemer wordt $3x + 3(-x) = 0$.
De limiet \textbf{bestaat niet} want de noemer is identiek $0$.\\

\textbf{Geval 2: $x \to +\infty$}\\

Voor $x > 0$ geldt $|x| = x$, dus:

$$\lim_{x \to +\infty} \dfrac{2 - |x|}{3x + 3|x|}
= \lim_{x \to +\infty} \dfrac{2 - x}{3x + 3x}
= \lim_{x \to +\infty} \dfrac{2 - x}{6x}
= \lim_{x \to +\infty} \dfrac{-x}{6x}
= -\dfrac{1}{6}$$

\begin{thebibliography}{999}
%\bibitem{a} Referentie
\end{thebibliography}
\end{document} 
